\subsection*{الف}
با قرار دادن $u=v=0$ داریم:
\begin{flalign*}
F(0,0) &= \frac{2 C(0) C(0)}{\sqrt{MN}}\sum_{i=0}^{7}\sum_{j=0}^{7} f(i,j)
= \frac{1}{\sqrt{64}}\sum_{i=0}^{7}\sum_{j=0}^{7} f(i,j) \\
&= \frac{1}{8}\times 5440 = 680 &&
\end{flalign*}

\subsection*{ب}
ضریب \lr{DC} همان مؤلفه فرکانس صفر است و بیانگر مقدار ثابت (میانگین روشنایی) تصویر است:
\begin{flalign*}
F(0,0) = \frac{1}{\sqrt{MN}}\sum f(i,j) = \sqrt{MN}\,\overline{f}, \quad \overline{f}=\frac{1}{MN}\sum f(i,j)
\end{flalign*}
در اینجا $\overline{f}=85$ و بنابراین $F(0,0)=8\times 85=680$؛ چون این مؤلفه متناظر با سیگنال ثابت (بدون نوسان) است، به آن \lr{DC} می‌گویند.

\subsection*{ج}
تابع پایه دوبعدی را تعریف می‌کنیم:
\begin{flalign*}
\phi_{u,v}(i,j)=\frac{2 C(u) C(v)}{\sqrt{MN}}
\cos\!\left(\frac{(2i+1)u\pi}{2M}\right)
\cos\!\left(\frac{(2j+1)v\pi}{2N}\right)
\end{flalign*}
پس
\begin{flalign*}
F(u,v)=\sum_{i=0}^{7}\sum_{j=0}^{7} f(i,j)\,\phi_{u,v}(i,j)
\end{flalign*}
برای $u=0,\ v=2$ داریم:
\begin{flalign*}
F(0,2)=\frac{2 C(0) C(2)}{\sqrt{MN}}\sum_{i=0}^{7}\sum_{j=0}^{7} f(i,j)
\cos\!\left(\frac{(2j+1)2\pi}{2N}\right)
\end{flalign*}
چون $\cos\!\left(\frac{(2i+1)0\pi}{2M}\right)=1$، الگو در راستای $i$ ثابت و فقط در راستای $j$ نوسانی با فرکانس $2$ دارد؛ این دقیقا همان الگوی تابع پایه در خانه $(0,2)$ در شکل است. بنابراین $F(0,2)$ میزان همبستگی تصویر با آن الگو را می‌سنجد.